\noindent\textbf{Elección del procedimiento.} La ley $\mathrm{Re}_c=C\mathrm{Ar}^n$ es el ajuste de datos docentes declarado, no una identidad universal. Se calcula Arquímedes y se evalúa la correlación; después se invierte $\mathrm{Re}=\rho Ud/\mu$ para obtener velocidad.

\begin{enumerate}

\item Usar el nuevo caso en el grupo de flotabilidad.

\[\fitmath{\mathrm{Ar}=\frac{\rho\cdot B\cdot {d}^{3}}{{\mu}^{2}}}\]

\[\fitmath{\mathrm{Ar}=\frac{(1000)\cdot (30)\cdot {(0.004)}^{3}}{{(0.001)}^{2}}\approx 1920\,}\]

\item Evaluar la potencia con el coeficiente y exponente del ajuste.

\[\fitmath{\mathrm{Rec}=C\cdot {\mathrm{Ar}}^{n}}\]

\[\fitmath{\mathrm{Rec}=(19.948)\cdot {(1920)}^{(0.30035)}\approx 193.2141\,}\]

\item Despejar velocidad de la definición de Reynolds.

\[\fitmath{\mathrm{Uc}=\frac{\mu\cdot \mathrm{Rec}}{\rho\cdot d}}\]

\[\fitmath{\mathrm{Uc}=\frac{(0.001)\cdot (193.2141299)}{(1000)\cdot (0.004)}\approx 0.04830353\,\mathrm{m/s}}\]

\end{enumerate}
